% chapters/ch07_dendrograms.tex
\chapter{Hierarchical Clustering as a Branching Object}
\label{ch:dendrograms}

\epigraph{%
  A dendrogram is a hypothesis about the history of similarity.%
}{\textit{---Flyxion}}

This chapter introduces dendrograms---the primary data structure of hierarchical clustering---and shows how they can be interpreted both combinatorially, as laminar families of subsets, and geometrically, as ultrametric spaces. The key object is the merge height: the value attached to each internal node of a dendrogram that records the scale at which two clusters merged. Reading the dendrogram at a fixed height threshold produces a flat clustering; varying the threshold produces the full hierarchy. We present two classical algorithms for constructing dendrograms---agglomerative (bottom-up) and divisive (top-down)---and note that neither optimizes any explicitly stated global objective. This motivates the formal objectives introduced in Chapter~\ref{ch:objectives} and the oracle-augmented methods of Part~IV. The visual comparison between a dendrogram and its corresponding ultrametric distance matrix---shown side by side in Figure~\ref{fig:dendrogram_matrix}---illustrates concretely the equivalence proved in Appendix~\ref{app:ultrametric_equiv}.

%% ── Figure: dendrogram with merge heights and matrix ─────────
\begin{figure}[ht]
\centering
% Dendrogram
\begin{tikzpicture}[scale=0.95]
% leaves at y=0
\foreach \x/\lbl in {0/A, 1/B, 2.3/C, 3.3/D} {
  \node[leafnode, label=below:{\small $\lbl$}] (\lbl) at (\x, 0) {};
}
% internal node AB at height 1
\node[internalnode, label=left:{\footnotesize $h=1$}] (AB) at (0.5, 1.2) {};
\draw[treeedge] (AB) -- (A);  \draw[treeedge] (AB) -- (B);

% internal node CD at height 1.5
\node[internalnode, label=right:{\footnotesize $h=1.5$}] (CD) at (2.8, 1.6) {};
\draw[treeedge] (CD) -- (C);  \draw[treeedge] (CD) -- (D);

% root at height 3
\node[internalnode, label=left:{\footnotesize $h=3$}] (R) at (1.65, 3.0) {};
\draw[treeedge] (R) -- (AB); \draw[treeedge] (R) -- (CD);

% height axis
\draw[dashed, nodegray!50] (-0.4, 1.2) -- (0.3, 1.2);
\draw[dashed, nodegray!50] (-0.4, 1.6) -- (0.3, 1.6);  % approximate
\draw[dashed, nodegray!50] (-0.4, 3.0) -- (0.3, 3.0);
\draw[->, >=stealth, nodegray!70] (-0.5, 0) -- (-0.5, 3.4)
  node[above, font=\tiny] {height};
\end{tikzpicture}
\hspace{1.8cm}
% Ultrametric distance matrix
\begin{tikzpicture}[scale=0.95]
\node[font=\small, above] at (2,3.4) {Ultrametric distance matrix $d_T$};
% Grid
\foreach \i/\lbl in {0/A, 1/B, 2/C, 3/D} {
  \node[font=\small] at (-0.4, 2.8-\i*0.75) {\lbl};
  \node[font=\small] at (0.6+\i*0.85, 3.1)  {\lbl};
}
% Fill cells
\foreach \row/\col/\val/\clr in {
  0/0/0/nodegray!5,    1/0/1/nodegreen!25, 2/0/3/nodeblue!20, 3/0/3/nodeblue!20,
  0/1/1/nodegreen!25, 1/1/0/nodegray!5, 2/1/3/nodeblue!20, 3/1/3/nodeblue!20,
  0/2/3/nodeblue!20, 1/2/3/nodeblue!20, 2/2/0/nodegray!5, 3/2/1.5/nodered!20,
  0/3/3/nodeblue!20, 1/3/3/nodeblue!20, 2/3/1.5/nodered!20, 3/3/0/nodegray!5
} {
  \fill[\clr] (0.2+\col*0.85, 2.5-\row*0.75) rectangle (0.95+\col*0.85, 3.1-\row*0.75);
  \draw[nodegray!40] (0.2+\col*0.85, 2.5-\row*0.75) rectangle (0.95+\col*0.85, 3.1-\row*0.75);
  \node[font=\footnotesize] at (0.575+\col*0.85, 2.8-\row*0.75) {\val};
}
\node[font=\footnotesize, nodegreen!80, below] at (0.9, -0.1) {merge at $h=1$};
\node[font=\footnotesize, nodered!80, below] at (2.9, -0.5) {merge at $h=1.5$};
\node[font=\footnotesize, nodeblue!80, below] at (2.0, -0.9) {merge at $h=3$};
\end{tikzpicture}
\caption{A dendrogram on four leaves (left) and its corresponding ultrametric distance matrix (right). Each internal node at height $h$ contributes the value $h$ to the distance between any pair of leaves whose LCA it is. The matrix encodes exactly the same information as the tree.}
\label{fig:dendrogram_matrix}
\end{figure}

\section{Dendrograms and Laminar Families}

\begin{definition}[Dendrogram]
A \emph{dendrogram} over a ground set $V$ is a rooted binary tree $T$ with $L(T) = V$ equipped with a node height function $h : V(T) \to \mathbb{R}_{\geq 0}$.
\end{definition}

The height function encodes the scale at which clusters merge. Reading the tree at threshold $\theta$ produces a flat clustering.

\begin{definition}[Horizontal cut]
For a dendrogram $T$ and threshold $\theta \geq 0$, the \emph{$\theta$-cut clustering} partitions $V$ into parts that are maximal subtrees whose root has height $\leq \theta$.
\end{definition}

\begin{proposition}
The collection of all parts of all $\theta$-cut clusterings forms a laminar family on $V$.
\end{proposition}

\section{Cluster Merges and the UPGMA Algorithm}

\begin{algorithm}[UPGMA / Average-linkage]
\label{alg:upgma}
\textbf{Input}: Similarity matrix $w_{ij}$.
\textbf{Initialize}: Clusters $\mathcal{C} = \{\{i\} : i \in V\}$.
\textbf{Repeat} until $|\mathcal{C}| = 1$: find $A, B \in \mathcal{C}$ maximizing average inter-cluster similarity, merge into $A \cup B$ at height $h = \bar{w}(A,B)$, update $\mathcal{C}$.
\textbf{Output}: the sequence of merge events defines a dendrogram $T$.
\end{algorithm}

UPGMA runs in $O(n^2 \log n)$ but does not optimize any global objective.

\section{Divisive Clustering}

Divisive methods begin with the full dataset and recursively partition it. The splitting oracle model (Part~IV) provides a principled way to perform the partition step using external information.

\section{Exercises}

\begin{enumerate}
  \item Given $w(1,2)=5$, $w(1,3)=3$, $w(2,3)=4$, apply UPGMA and draw the resulting dendrogram.
  \item Show that single and complete linkage can produce different dendrograms on the same similarity matrix by exhibiting an example on four points.
  \item From the dendrogram in Figure~\ref{fig:dendrogram_matrix}, read off the flat clustering at threshold $\theta = 2$.
\end{enumerate}
